\documentclass[10pt,a4paper]{amsart}
\usepackage[margin=19mm]{geometry}
\usepackage{amsmath,amssymb,mathtools}
\usepackage[scr=rsfs,cal=euler]{mathalfa}
\usepackage{xfrac}
\usepackage[unicode,hidelinks,bookmarks=false]{hyperref}
\newcommand{\ie}{i.e.\ }
\newcommand{\cf}{cf.\ }
\newcommand{\resp}{respectively\ }
\newcommand{\Subsection}{\subsection}
\providecommand{\nobreakdash}{\nobreak\hbox{-}}
\setlength{\emergencystretch}{2em}
\multlinegap0pt
\usepackage{mdframed}
\usepackage{mdframed}
\usepackage{mdframed}
\usepackage{mdframed}
\usepackage{tikz}
\usepackage{amssymb}
\usepackage{float}
\usepackage{xcolor}
\usepackage{tcolorbox}
\usepackage{mdframed}
\usepackage{algorithm}\usepackage{algpseudocode}
\newtheorem{proposition}{Proposition}
\title{Sharp Support Thresholds for Smeariness of Absolutely Continuous Measures on Spheres}
\author{Susovan Pal}
\date{Source: arXiv:2606.02144v1 (CC BY 4.0)}
\begin{document}
\maketitle

\noindent\emph{Source and licence.} Sharp Support Thresholds for Smeariness of Absolutely Continuous Measures on Spheres. Version arXiv:2606.02144v1 (CC BY 4.0), distributed by arXiv under the Creative Commons Attribution 4.0 International licence, http://creativecommons.org/licenses/by/4.0/, as recorded in the arXiv licence metadata for this exact version and shown on the abstract page https://arxiv.org/abs/2606.02144v1. Full licence notice: \texttt{licenses/arxiv-2606.02144-CC-BY-4.0.txt}.\par\smallskip

\noindent\textit{Editorial scope.} Counted unit: the article's proposition prop:asymptotics-Rm-Sm (source0.tex lines 1705--1715). The article's complete original proof of it is delivered with it (source0.tex lines 1717--1795, one \texttt{proof} environment of 79 lines). No mathematics is written, repaired or re-derived: the statement and the proof are the article's own lines, in the article's own notation, and no retained line is substituted or re-flowed.  Retained uncounted as the source-local context the proof needs: lines 473--489; lines 688--690; lines 739--743; lines 785--798.  Nothing was omitted from any retained span, so the package carries no omission record.  No retained line contains a citation, so the wrapper carries no bibliography and no bibliography entry.  No retained line contains a figure environment, an image-inclusion command or drawing code, so no image is delivered and the package declares no asset dependency.  Editorial formatting only, in addition: an AMS \texttt{amsart} preamble that restates the article's own declarations and macros used by the retained lines, namely the environments \texttt{proposition} on separate counters, together with the standard packages those lines need.  The retained environments print in this document as Proposition 1 in order of appearance, whereas the article prints the same items with the numbering of its own full document.  The complete native source of the article as distributed in the arXiv e-print for this exact version is bundled unchanged as \texttt{sources/201881/source0.tex}.
\begin{equation}\label{eqn:A-B}
\boxed{
\begin{aligned}
A_{3}(R)&= - \frac{2}{3} + \frac{1}{\sin^{2} R} - \frac{1}{R \tan^{5} R}
- \frac{2\cos R}{R \sin^{3} R} + \frac{\cos R}{R \sin^{5} R},
\qquad
B_{3}(R)= \frac{2}{3 R^{2}} - \frac{1}{R^{2}\sin^{2} R} + \frac{1}{R^{3}\tan R},\\
A_{4}(R)&=\frac{\frac{6 R}{\tan R} - \frac{15 R \cos R}{\sin^{3} R} - 11 + \frac{15}{\sin^{2} R}}{12 R^{4}},\\
A_{2}(R)&=\frac{- \frac{4 R}{\tan^{5} R} + \frac{R \cos R}{\sin^{3} R} + \frac{4 R \cos R}{\sin^{5} R}
- \frac{7}{\tan^{6} R} + \frac{12}{\sin^{2} R} - \frac{21}{\sin^{4} R} + \frac{7}{\sin^{6} R}}{6 R^{2}},\\
A_{0}(R)&=
- \frac{R}{6 \tan^{7} R} + \frac{R \cos^{5} R}{4 \sin^{7} R} - \frac{R \cos R}{12 \sin^{7} R}
- \frac{1}{4 \tan^{8} R} - \frac{3}{4 \sin^{4} R} + \frac{5}{4 \sin^{6} R}
+ \frac{3 \cos^{6} R}{4 \sin^{8} R} - \frac{1}{2 \sin^{8} R}.
\end{aligned}
}
\end{equation}

\begin{equation}\label{eqn:b}
    b_m(R):=1+(m-1)R\cot R .
\end{equation}

\begin{equation}
\label{eqn:h}
    h_m(R):=
    A_0(R)+\frac{1}{m}A_2(R)R^2+\frac{3}{m(m+2)}A_4(R)R^4 .
\end{equation}

\begin{proposition}[\textbf{The Hessian threshold}]
\label{prop:zero_bm}
Let \(m\ge2\), and following \eqref{eqn:b}, let \(b_m:(0,\pi)\to\mathbb R\) be defined by
\[
    b_m(R)=1+(m-1)R\cot R.
\]
Then \(b_m\) is strictly decreasing on \((0,\pi)\), has a unique zero
\begin{equation}\label{eqn:R_m}
    \boxed{R_m\in(\pi/2,\pi),}
\end{equation}
    

and satisfies \(b_m(R)>0\) for \(R\in(0,R_m)\) and \(b_m(R)<0\) for \(R\in(R_m,\pi)\).
\end{proposition}

\begin{proposition}[\textbf{High-dimensional location of \(R_m\) and \(S_m\)}]
\label{prop:asymptotics-Rm-Sm}
\label{prop:asymp-first-zero-h_m}
As \(m\to\infty\),
\[
R_m=\frac{\pi}{2}+O\!\left(\frac{1}{m}\right),
\qquad
S_m=\frac{\pi}{2}+O\!\left(\frac{1}{m}\right).
\]
In particular, \(R_m\to\pi/2\) and \(S_m\to\pi/2\).
\end{proposition}

\begin{proof}
\textbf{Analysis of \(R_m\).} By Proposition~\ref{prop:zero_bm}, \(R_m\) is the unique zero of \(b_m(R)=1+(m-1)R\cot R\), and \(R_m\in(\pi/2,\pi)\). Write
\[
R_m=\frac{\pi}{2}+\delta_m,
\qquad
n:=m-1,
\]
so that \(\delta_m\in(0,\pi/2)\). Since \(\cot(\pi/2+\delta)=-\tan\delta\), the equation \(b_m(R_m)=0\) becomes
\[
\left(\frac{\pi}{2}+\delta_m\right)\tan\delta_m=\frac{1}{n}.
\]
As \(0<\delta_m<\pi/2\), we have \(\delta_m\le\tan\delta_m\). Hence
\[
\frac{\pi}{2}\delta_m
\le
\left(\frac{\pi}{2}+\delta_m\right)\tan\delta_m
=
\frac{1}{n}.
\]
Therefore \(0<\delta_m\le 2/(\pi n)\), and hence
\[
R_m=\frac{\pi}{2}+O\!\left(\frac{1}{m}\right).
\]

\textbf{Analysis of \(S_m\).}  Set \(S_0:=\pi/2\). From \eqref{eqn:h},
\[
h_m(R)=A_0(R)+\frac{1}{m}A_2(R)R^2+\frac{3}{m(m+2)}A_4(R)R^4.
\]
Direct evaluation of the coefficients in \eqref{eqn:A-B} at \(S_0\) gives
\[
A_0(S_0)=0,
\qquad
A_0'(S_0)=\frac{\pi}{24}>0,
\qquad
A_2(S_0)S_0^2=-\frac{1}{3}.
\]
Also \(A_4(S_0)S_0^4\) is finite. Therefore
\[
h_m(S_0)
=
-\frac{1}{3m}
+
O\!\left(\frac{1}{m^2}\right)
<0
\]
for all sufficiently large \(m\).

Since \(A_0'(S_0)=\pi/24\), there exists \(r>0\) such that
\[
A_0(S_0+\delta)\ge \frac{\pi}{48}\delta
\qquad
\text{for all }0\le\delta\le r.
\]
Moreover, \(A_2(R)R^2\) and \(A_4(R)R^4\) are bounded on \([S_0,S_0+r]\). Hence there exists \(C>0\) such that, for all \(0\le\delta\le r\) and all \(m\ge4\),
\[
h_m(S_0+\delta)
\ge
\frac{\pi}{48}\delta-\frac{C}{m}.
\]
Choose \(\Delta>48C/\pi\). For all sufficiently large \(m\), we have \(\Delta/m\le r\), and therefore
\[
h_m\!\left(S_0+\frac{\Delta}{m}\right)
\ge
\frac{1}{m}\left(\frac{\pi\Delta}{48}-C\right)>0.
\]
Since \(h_m(S_0)<0\) and \(h_m\) is continuous, \(h_m\) has a zero in
\[
\left(S_0,S_0+\frac{\Delta}{m}\right).
\]
By definition, \(S_m\) is the first zero of \(h_m\) in \((\pi/2,\pi)\). Consequently,
\[
0<S_m-S_0\le \frac{\Delta}{m}
\]
for all sufficiently large \(m\). Thus
\[
S_m=\frac{\pi}{2}+O\!\left(\frac{1}{m}\right),
\]
and the proof is complete.
\end{proof}

\end{document}
